\section{Broader Lessons and Limitations} \label{section:broader_lessons_and_limitations}

MLPs appear reliable for ranking interfacial candidates when the bonding environments and geometries fall within the training domain. Screening workflows benefit most from MLPs in systems involving Ge, Sn, and well-registered C|X structures, where rank preservation is consistently strong. In contrast, strained or covalently dominated systems like Si|Ge or Si|C exhibit erratic performance, especially when interfacial bonding is asymmetric or poorly coordinated.

These results highlight the need for delta-correction strategies or domain-specific retraining, particularly for silicon-rich or mixed-interface systems. Hybrid workflows that combine MLP-based pre-screening with selective DFT confirmation offer a scalable and pragmatic approach to large-scale interface discovery.

This study reinforces the value of rank-based benchmarking — such as Spearman, Kendall, and Top-N overlap — over absolute energy metrics. In screening contexts, the ability to reliably identify top-performing configurations is more actionable than minimising raw energy error.